% native creation stories: indigenous-b

\subsection{Slavic}

At the root is \textbf{Rod}, the undivided source of kin and birth, with the \textbf{Rozhanitsy}, the spirits of birth and fate. At the foundation stand a dual pair, \textbf{Belobog}, the White God of light and order, and \textbf{Chernobog}, the Black God of darkness and disorder. In the beginning there is only the primordial sea and a creating power with a diver, a bird or the devil. The diver goes to the depths and brings up a handful of mud. The creating power scatters it on the water and it grows into the ordered land. The diver withholds part of the mud by deceit, and the hidden portion swells and is cast out to become the mountains and swamps. Over the world rises the \textbf{world tree}, an oak, with roots in the depths, trunk among the living, and crown in the heavens. Across sky and deep \textbf{Perun} the thunder god and \textbf{Veles} the underworld god contend forever.

\subsection{Baltic}

At the head stands \textbf{Dievas}, the sky father, who fixes the order of the heavens. From him come the divine twins, the sons of the sky. The visible world is governed by a celestial family. \textbf{Saule}, the sun, is the source of all life and carries the souls of the dead across the sea, dying and reborn each day. \textbf{Menulis} the moon, in Latvian \textbf{Meness}, is set against her as her partner. The moon is unfaithful with the morning star \textbf{Ausrine}, and \textbf{Perkunas} the thunder god, in Latvian \textbf{Perkons}, splits him, so day and night are divided. On the path of the sun, by the water at the edge of the world, stands the \textbf{Austras koks}, the tree of dawn, joining the realms of the sky. Around it turns the recurring celestial wedding, where the \textbf{Saules meita}, the daughter of the sun, is courted by the sons of the sky, the morning star, and the thunderer.

\subsection{Shinto}

The world is full of \textbf{kami}, the powers of things. The first powers come into being alone in the high plain of heaven and conceal themselves, carrying \textbf{musubi}, the force of growth, in their names. After seven generations come the primal pair, \textbf{Izanagi} and \textbf{Izanami}. Standing on the floating bridge of heaven, they lower a jeweled spear into the brine and stir it, and the falling drops pile up into the first island. Descending, they wed and bring forth the islands and many kami. Bearing the fire power, Izanami is burned and dies into the land of the dead. Izanagi follows, flees her corrupted form, and cleanses himself in a river. From the cleansing come three great powers, \textbf{Amaterasu} the sun from one eye, \textbf{Tsukuyomi} the moon from the other, and \textbf{Susanoo} the storm from the nose. When the storm runs wild the sun hides in the rock cave, the \textbf{Ama-no-Iwato}, and the world goes dark until the powers draw her out with a mirror and a dance. At last she sends her grandchild \textbf{Ninigi} down with the mirror, sword, and jewel, the \textbf{tenson korin} that founds the ruling line.

\subsection{Tamil}

The one reality is at once the still ground, the unmoving witness, and the moving power, the dance of \textbf{Shiva} and \textbf{Shakti} that brings out the worlds. These are not two beings but one seen as rest and motion. The motion plays out as five acts performed forever. The power brings the worlds into being, holds them, draws them back, veils itself so the worlds seem separate, then shows itself again and frees the bound soul. The same presence that fills the high heaven fills the land, the body, and the breath. To find it the soul turns inward and downward into its own body, held to be a temple of the whole. The path binds yoga, the inner fire of \textbf{kundalini}, and the worship of Shiva, as taught in the \textbf{Tirumantiram} of \textbf{Tirumular}. The \textbf{Siddhars} sought the eight \textbf{ashta siddhi} and a deathless body through \textbf{kaya kalpa}, and the older \textbf{Sangam} poetry mapped the divine onto the land through the five \textbf{tinai}.

\subsection{Polynesian}

In the \textbf{Kumulipo} the world begins in \textbf{po}, the deep and slimy darkness that is full of latent life. From it the first living thing arises, the coral polyp, and after it come in ordered stages the creatures of the sea, then the land things, then the gods and humankind, a slow climb from the dark toward \textbf{ao}, the light. In the Maori telling the movement runs through three states, \textbf{Te Kore} the void, \textbf{Te Po} the long night, and \textbf{Te Ao} the world of light. The sky father \textbf{Ranginui} and the earth mother \textbf{Papatuanuku} lie locked in a tight embrace, and their children are crushed in the narrow dark between them. The children take counsel, and one after another they try to force the parents apart and fail, until \textbf{Tane}, the forest god, braces his back against the earth and his feet against the sky and pushes them open. Light floods in and the living world begins. \textbf{Tawhirimatea}, the storm god, will not accept the parting and rages against his brothers forever.

\subsection{Aboriginal Australian}

In the formative time, the \textbf{Dreaming}, the land lay unshaped and asleep, a featureless plain. The ancestral beings rose from beneath the ground or came across the empty country. As they walked, hunted, fought, and rested, the hills, rivers, waterholes, and rocks took their present forms. They made the plants and animals, set down the law, and gave each group its land and its language. Among them is the \textbf{Rainbow Serpent}, the water-creator known by many names such as \textbf{Ngalyod}, \textbf{Wagyl}, \textbf{Yurlunggur}, and \textbf{Wanampi}, who shaped deep channels and ridges and who both gives life and brings danger. The paths the ancestral beings traveled are remembered as \textbf{songlines}, sung tracks that record the journeys and let people read and cross the land. When their work was done some beings sank back into the ground or rose into the sky. The formative time never closed. It remains present in the land, in ritual, and in the law.

\subsection{Korean}

In the \textbf{Cheonjiwang Bonpuri} of the Jeju shamans, heaven and earth are at first one fused mass with no gap. Then the mass splits. The light rises to become the sky and the heavy settles to become the ground, and the \textbf{gawp} between them opens the room of the world. But the early world is unlivable, for two suns blaze by day and two moons freeze the night. The heavenly king \textbf{Cheonjiwang} comes down and fathers twin sons by the earthly woman \textbf{Bagi}, the Great Star King \textbf{Daebyeol-wang} and the Little Star King \textbf{Sobyeol-wang}. The brothers climb to the sky and take up iron bows, and one shoots down the extra sun and the other the extra moon, leaving a single sun and a single moon. Then they hold a flower-growing contest to decide who rules the living and who the dead. The younger Sobyeol-wang wins the world of the living by a trick, which is why it still holds disorder, while the elder Daebyeol-wang takes the world of the dead and rules it with strict justice. In the older mainland strand the grandmother giant \textbf{Mago}, also called \textbf{Magohalmi}, shapes the mountains and rivers by her own body.

\subsection{Yoruba (Ifa)}

In the beginning there is only the sky above, ruled by \textbf{Olodumare} the supreme creator, and a wide marsh of water below, with no solid ground. Olodumare gives \textbf{Obatala}, guided by \textbf{Orunmila} the power of wisdom, the task of making land. Obatala comes down on a chain carrying a snail shell full of earth, a hen, and a pigeon. He pours out the earth on the water and sets the birds down, and they scratch and spread it in every direction until it becomes the wide dry land. The first sacred city, Ile-Ife, is founded on that spot. Olodumare then breathes life into the world and sets the many \textbf{orisha} to govern its parts, their work driven by \textbf{ashe}, the life-force that makes things come to pass. Each person, before birth, kneels before the creator and chooses an inner head, the \textbf{ori}, that carries a destiny. The work of a life is to align with it. The teaching is preserved and read through the \textbf{Ifa} system, whose corpus is organized into the \textbf{odu}, the fixed set of signs.

\subsection{Inuit}

The ordering principle of the world is \textbf{Sila}, the breath and the air. It is the invisible force that fills everything, governs the weather and the winds, and gives each living thing the breath of its life. It is at once the air outside and the breath inside, and it holds the moral order, so that breaking the rules of right living disturbs it and the weather turns harsh. The food of the sea comes from \textbf{Sedna}, the woman at the bottom of the ocean. She was thrown into the water and her fingers were cut away as she clung to the boat, and from the severed fingers came the seals, the walruses, and the whales. She rules \textbf{Adlivun}, the sea-floor underworld, and holds the animals. When people break the rules she grows angry, tangles her hair, and keeps the animals down, and the people starve. Then the \textbf{angakkuq}, the spirit-specialist, travels to her, calms her, and combs her hair so she releases the animals again. The raven \textbf{Tulugaq}, creator and trickster, takes the form of a bird, brings light, and stirs the storms, and the sun and the moon are siblings who fled each other into the sky.
